\section[$B^{+}_{\mathrm{dR}}$-仿射 Grassmann 簇]{$B^{+}_{\mathrm{dR}}$-仿射 Grassmann 簇（12 月 7 日）}

\srcpage{60}

\subsection{回顾}

固定通常的记号：$E\supset\cO_E\ni\pi$，剩余域为 $\Fq$。取定 $\ol{\Fq}$，得到 $\breveE=W_{\cO_E}(\ol{\Fq})[1/\pi]$。

设 $G/E$ 是约化群。上次我们讨论了「带 $G$-结构的等晶体（isocrystal）」这一概念，其同构类记作 $B(G)$。这是一个可数集，可以具体表示成 $G(\breveE)/\sigma$-共轭，即 $b\in G(\breveE)$ 模去关系 $b\sim g^{-1}b\sigma(g)$。

我们定义了两个映射
\[
\nu\colon B(G)\to (X^{+}_{\QQ})^{\Gamma},\qquad \Gamma=\Gal(\ol E/E),
\]
以及
\[
\kappa\colon B(G)\to \pi_1(G)_{\Gamma}.
\]
要点在于 $(\nu,\kappa)\colon B(G)\to (X^{+}_{\QQ})^{\Gamma}\times\pi_1(G)_{\Gamma}$ 是单射，因此这用「组合」数据分类了 $B(G)$。

\subsection{$\Bun_G$ 的结构}

\begin{definition}
$\Bun_G$ 是 $\Perf_{\Fq}$ 上的 $v$-叠，把 $S$ 送到 $X_S$ 上 $G$-主丛的群胚。
\end{definition}

注意，对任意 $S$ 都有函子 $G\text{-}\Isoc\to\Bun_G(S)$，它可以看作沿 $X_S=Y_S/\varphi^{\ZZ}\to\Spa\breveE/\sigma^{\ZZ}$ 的拉回。

\begin{remark}
事实上，$G$-等晶体恰好就是「$X_S$ 上的泛 $G$-主丛」，这由下面的定理给出。
\end{remark}

\begin{theorem}[Anschütz]
函子
\[
G\text{-}\Isoc\to\varprojlim_{S}\Bun_G(S)
\]
是等价。
\end{theorem}

\begin{theorem}[$E/\Qp$ 时由 Fargues，一般情形由 Anschütz]
若 $S=\Spa(C,C^+)$，其中 $C$ 是完全代数闭域，则函子
\[
G\text{-}\Isoc\to\Bun_G(S)
\]
把 $\Spa\breveE/\sigma^{\ZZ}$ 上的 $G$-主丛送到其沿 $Y_S/\varphi^{\ZZ}=X_S$ 的拉回，并在同构类上诱导双射
\[
\Bun_G(S)/\!\sim\ \xrightarrow{\ \sim\ } B(G).
\]
\end{theorem}

\begin{theorem}
映射 $|\Bun_G|\to B(G)$ 连续，即
\begin{enumerate}[label=(\arabic*)]
\item $-\nu\colon|\Bun_G|\to(X^{+}_{\QQ})^{\Gamma}$ 半连续；
\item $-\kappa\colon|\Bun_G|\to\pi_1(G)_{\Gamma}$ 局部常值。
\end{enumerate}
事实上，$\kappa$ 诱导双射
\[
\kappa\colon\pi_0(\Bun_G)\xrightarrow{\sim}\pi_1(G)_{\Gamma}.
\]
\end{theorem}

\begin{remark}
$|\Bun_G|$ 的完全确定由 Hansen 对 $\mathrm{GL}_n$ 得到，由 Hamann 对一些经典群得到，由 Viehmann 对一般 $G$ 得到。
\end{remark}

\begin{proof}
(1) 我们对 $\mathrm{GL}_n$ 已知。Rapoport–Richartz [RR96] 解释的一个简单论证表明，可以通过考虑嵌入 $G\hookrightarrow\mathrm{GL}_n$ 化归到 $\mathrm{GL}_n$。

(2) 这更难；事实上，Kottwitz 映射 $\kappa$ 的构造本身就很复杂！证明沿用构造中同样的群论手法。
\end{proof}

\srcpage{61}

\begin{lemma}
设 $G'\twoheadrightarrow G$ 是群之间的映射，且是以中心环面为核的扩张（即 $z$-扩张）。则诱导映射 $\Bun_{G'}\to\Bun_G$ 是 $v$-叠的满射。
\end{lemma}

我们稍后给出证明。经典情形下它归结为：提升的障碍是一个 $H^2$，而它在曲线上消失。

暂设该引理，我们来完成上述定理的证明。步骤与 $\kappa$ 的构造平行：

\begin{proof}[上述定理的证明（续）]
\begin{itemize}
\item 若 $G$ 是诱导环面（induced torus）的乘积，则 $\pi_1(G)_{\Gamma}$ 无挠。
\[
\begin{tikzcd}
& (X^{+}_{\QQ})^{\Gamma} \arrow[dr] & \\
B(G) \arrow[ur, "\nu"] \arrow[dr, "\kappa"'] & & (\pi_1(G)_{\QQ})^{\Gamma} \\
& \pi_1(G)_{\Gamma} \arrow[ur, "\text{平均}"'] &
\end{tikzcd}
\]
由于 $\pi_1(G)_{\Gamma}$ 无挠，此时 $\kappa$ 由 $\nu$ 决定。$(X^{+}_{\QQ})^{\Gamma}$ 中唯一的非平凡序关系位于到 $(\pi_1(G)_{\QQ})^{\Gamma}$ 的纤维中，因此 $\kappa$ 局部常值。
\item 对一般的环面 $T$，可取满射 $\wt T\to T$，其中 $\wt T$ 是诱导环面的乘积。由上述引理，
\[
\Bun_{\wt T}\to\Bun_T
\]
是 $v$-叠的满射。特别地，$|\Bun_{\wt T}|\to|\Bun_T|$ 是商映射\footnote{见 [S17, Proposition 11.13]。}。所以 $\kappa_T\colon|\Bun_T|\to\pi_1(T)_{\Gamma}$ 的连续性由刚刚建立的 $\kappa_{\wt T}\colon|\Bun_{\wt T}|\to\pi_1(\wt T)_{\Gamma}$ 的连续性推出。
\[
\begin{tikzcd}
\lvert\Bun_{\wt T}\rvert \arrow[r] \arrow[d, "\kappa_{\wt T}"'] & \lvert\Bun_T\rvert \arrow[d, "\kappa_T"] \\
\pi_1(\wt T)_{\Gamma} \arrow[r] & \pi_1(T)_{\Gamma}
\end{tikzcd}
\]
\item 考虑 $G_{\mathrm{der}}$ 单连通的 $G$。则有短正合列
\[
1\to G_{\mathrm{der}}\to G\to T\to 1,
\]
诱导 $\pi_1(G)\xrightarrow{\sim}\pi_1(T)$。于是有图
\[
\begin{tikzcd}
\lvert\Bun_G\rvert \arrow[r] \arrow[d, "\kappa_G"'] & \lvert\Bun_T\rvert \arrow[d, "\kappa_T"] \\
\pi_1(G)_{\Gamma} \arrow[r, "\sim"] & \pi_1(T)_{\Gamma}
\end{tikzcd}
\]
由于 $\kappa_T$ 连续，这也说明 $\kappa_G$ 连续。
\item 对一般的 $G$，取 $z$-扩张 $G'\to G$ 使 $G'_{\mathrm{der}}$ 单连通，然后应用上述引理。
\end{itemize}
\end{proof}

\srcpage{62}

\begin{corollary}[对 {[RR96]} 的略微加强]
若 $S$ 是 $\ol{\Fq}$ 上的完满概形，$\cG$ 是 $S$ 上的 $G$-等晶体，则
\[
\kappa\colon |S|\to\pi_1(G)_{\Gamma}
\]
局部常值。
\end{corollary}

\begin{proof}
不妨设 $S=\Spec R$。对任意带映射 $S'\to\Spa(R,R)$ 的完满空间 $S'/\ol{\Fq}$，用 $\cG$ 按同样的步骤得到映射 $S'\to\Bun_G$。对这个映射我们得到局部常值性结论。（不过需要小心：如我们所讨论的，主丛与等晶体的特殊化关系是反过来的。）
\end{proof}

\begin{remark}
Rapoport–Richartz 所缺的技术要素正是上述引理的某种类似物。
\end{remark}

所以现在回到上述引理的证明。

我们将从 Beauville–Laszlo 单值化的一个版本推出它。我们将引入 $B^+_{\mathrm{dR}}$-仿射 Grassmann 簇以及 $v$-叠的满射 $\Gr^{B^+_{\mathrm{dR}}}_G\to\Bun_G$。我们将看到，若 $G'\twoheadrightarrow G$ 是中心同源，则 $\Gr^{B^+_{\mathrm{dR}}}_{G'}\to\Gr^{B^+_{\mathrm{dR}}}_G$ 是满射。上述引理由此立即推出。

这个单值化是 §9.2 中映射 $\PP^1\to\Bun_2$ 的推广版本。后者从平凡丛出发，在一点处对它作修正而得到。一般情形的性质相同。

\subsection{$B^+_{\mathrm{dR}}$}

设 $R$ 是 $\cO_E$ 上任一完满环（perfectoid ring）。设 $\varpi\in R^{\flat}$ 是伪一致化子（pseudo-uniformizer）。则得到 Fontaine 映射
\[
\theta\colon W_{\cO_E}(R^{\flat\circ})\to R,
\]
其核为 $(\xi)$，其中 $\xi=\pi+[\varpi^{\epsilon}]a$，$a\in W_{\cO_E}(R^{\flat\circ})$。\yizhu{原文此处写作 $\varpi^{\epsilon}$，但 $\epsilon$ 在文中未定义；图版如此，照录。}{A（原文 p.62）}

所以由 $\theta$ 得到满射
\[
W_{\cO_E}(R^{\flat\circ})[1/[\varpi]]\twoheadrightarrow R=R^{\circ}[1/\varpi^{\#}].
\]

\begin{definition}
定义 $B^+_{\mathrm{dR}}(R)$ 为 $W_{\cO_E}(R^{\flat\circ})[1/[\varpi]]$ 的 $\xi$-进完备化。
\end{definition}

这在几何上意味着什么？

\begin{itemize}
\item 这是 $R$ 的「单参数形变」，特别地有
\[
B^+_{\mathrm{dR}}(R)/(\xi)\xrightarrow{\sim}R.
\]
特别地，$(\xi)^i/(\xi)^{i+1}$ 作为 $R$-模也同构于 $R$。
\item $B^+_{\mathrm{dR}}(R)\twoheadrightarrow R$ 是固体 $\cO_E$-代数中的泛 pro-无穷小加厚（pro-infinitesimal thickening）。
\item 若 $R/\Fq$，则 $B^+_{\mathrm{dR}}(R)=W_{\cO_E}(R)$。
\item 若 $R=C/E$ 完全代数闭，则 $B^+_{\mathrm{dR}}(R)$ 作为抽象 $E$-代数同构于 $C[[\xi]]$（但不同构为拓扑 $E$-代数——事实上 $B^+_{\mathrm{dR}}(R)$ 不是拓扑环，尽管它是带拓扑的环）。
\item 这个名字来自 Fontaine，它作为「$p$ 进 de Rham 周期环」出现。
\end{itemize}

\srcpage{63}
\begin{center}
\figP{063}
\end{center}

\subsection{$B^+_{\mathrm{dR}}$-仿射 Grassmann 簇}

\begin{definition}
$B^+_{\mathrm{dR}}$-仿射 Grassmann 簇 $\Gr^{B^+_{\mathrm{dR}}}_G$ 是 $\Perf/(\Spa E)^{\diamond}\cong\Perf/E$ 上函子的 étale 层化，该函子把 $\Spa(R,R^+)/E$ 送到 $G(\BdR(R))/G(B^+_{\mathrm{dR}}(R))$，其中 $\BdR(R)=B^+_{\mathrm{dR}}(R)[1/\xi]$。
\end{definition}

等价地，它分类 $\Spa B^+_{\mathrm{dR}}(R)$ 上的 $G$-主丛，连同在 $\Spa \BdR(R)$ 上的平凡化。原因是在 $R$ 上 étale 局部地，$\Spa B^+_{\mathrm{dR}}(R)$ 上的 $G$-主丛平凡。

\begin{remark}
这是通常仿射 Grassmann 簇的 $p$ 进版本
\[
R/\CC\mapsto G(R((t)))/G(R[[t]]).
\]
\end{remark}

\begin{example}
注意当 $E=\Fq((t))$ 时，这就是字面意义上的
\[
R/\Fq((t))\mapsto G(R((t-\zeta)))/G(R[[t-\zeta]]).
\]
这是因为若 $E=\Fq((t))$，则 $W_{\cO_E}(R^{\circ})=R^{\circ}[[t]]$，且 $\xi=t-\zeta$。因此这恢复了通常的仿射 Grassmann 簇。
\end{example}

\begin{proposition}[Cartan 分解]
我们有
\[
\Gr^{B^+_{\mathrm{dR}}}_G(C)=\bigcup_{\mu\in X^{+}}G(B^+_{\mathrm{dR}}(C))[\mu(\xi)].
\]
\end{proposition}

\begin{example}
对 $G=\mathrm{GL}_n$，若 $\mu=(a_1\ge\dots\ge a_n)$，则 $\mu(\xi)$ 为
\[
\begin{pmatrix}
\xi^{a_1} & & & \\
& \xi^{a_2} & & \\
& & \ddots & \\
& & & \xi^{a_n}
\end{pmatrix}.
\]
\end{example}

\begin{proof}
作为抽象环有同构 $B^+_{\mathrm{dR}}(C)\cong C[[\xi]]$。这里的问题关乎代数环，因为它可以归约为向量丛的问题，而后者只依赖底层的代数环。所以它实际上只是通常 Cartan 分解的推论。
\end{proof}

为简单起见，设 $G$ 分裂。（一般地，在 $E$ 上作有限 étale 基变换后如此，且 $\Gr_{G_{E'}}=\Gr_G\times_E E'$，故可用下降化归到此情形。）

\srcpage{64}

\begin{definition}
对 $\mu\in X^{+}$，令
\[
\Gr^{B^+_{\mathrm{dR}}}_{G,\le\mu}\subset\Gr^{B^+_{\mathrm{dR}}}_G
\]
为「Schubert 簇」子函子，由所有 $S\to\Gr^{B^+_{\mathrm{dR}}}_G$ 组成，使得在每个几何点处，它落在 $\mu'(\xi)$ 的 $G(B^+_{\mathrm{dR}})$-轨道中，其中 $\mu'\le\mu$。
\end{definition}

\begin{remark}
由于完满空间自动既约，可以通过在几何点上指定条件来指定子集。对概形来说这不是个好主意。
\end{remark}

\begin{theorem}
$\Gr^{B^+_{\mathrm{dR}}}_{G,\le\mu}\to(\Spa E)^{\diamond}$ 紧合，且 $\Gr^{B^+_{\mathrm{dR}}}_{G,\le\mu}$ 由空间菱形（spatial diamond）表示。
\[
\Gr^{B^+_{\mathrm{dR}}}_G=\bigcup_{\mu}\Gr^{B^+_{\mathrm{dR}}}_{G,\le\mu}
\]
且过渡映射是闭浸入。

若 $\mu$ 是极小的（minuscule），则 $\Gr^{B^+_{\mathrm{dR}}}_{G,\le\mu}\cong(G/P_\mu)^{\diamond}$。
\end{theorem}

证明有点微妙。它用到 Artin 对代数空间的识别原理的一个变体。在此情形很难找到显式的 pro-étale 覆盖。我们改为在 $v$-拓扑中找图册，这已足以表明它是空间的。然后有一个判据，用几何点刻画空间 $v$-层何时是菱形。几何点被各分层穷尽，而在分层上可以直截了当地论证。

\begin{corollary}
若 $G'\to G$ 是 $z$-扩张，则 $\Gr^{B^+_{\mathrm{dR}}}_{G'}\to\Gr^{B^+_{\mathrm{dR}}}_G$ 是 $v$-覆盖。
\end{corollary}

\begin{proof}
作域扩张，可设 $G$ 分裂。只需证明对一切 $\mu'\mapsto\mu$，
\[
\Gr^{B^+_{\mathrm{dR}}}_{G',\le\mu'}\to\Gr^{B^+_{\mathrm{dR}}}_{G,\le\mu}
\]
是 $v$-覆盖。这里用到 $(X')^{+}\twoheadrightarrow X^{+}$，因为 $G'\to G$ 的核是连通环面。这些都是空间菱形，特别地 qcqs。因此可在几何点上检验满射性。于是由 Cartan 分解与 $G'(B^+_{\mathrm{dR}}(C))\twoheadrightarrow G(B^+_{\mathrm{dR}}(C))$ 立得。
\end{proof}

\begin{remark}
对 $E/\Fq((t))$，所有 $\Gr_{G,\le\mu}$ 都由射影簇表示；但对 $E/\Qp$，我们得到真正的菱形，它们不来自刚性簇。
\end{remark}

\subsection{Beauville–Laszlo 单值化}

\begin{definition}
设 $G$ 是 $E$ 上的约化群（不必分裂）。有映射
\[
\Gr^{B^+_{\mathrm{dR}}}_G\to\Bun_G,
\]
定义如下：定义域的一个点是 $S=\Spa(R,R^+)/E$，以及 $B^+_{\mathrm{dR}}(R)$ 上的 $G$-主丛 $\cE_0$，它在 $\BdR(R)$ 上平凡化。我们可以把 $X_S^{\mathrm{alg}}\setminus\Spec R$ 上的平凡 $G$-主丛与 $\cE_0/B^+_{\mathrm{dR}}(R)$ 沿 $\BdR(R)$ 上给定的等同粘合起来。
\end{definition}

\srcpage{65}
\begin{center}
\figP{065}
\end{center}

我们在这里用到：

\begin{lemma}[Beauville–Laszlo]
若 $X$ 是 $E$ 上的概形，$Z\subset X$ 是 Cartier 除子且 $Z$ 仿射，则 $X$ 上的 $G$-主丛等价于下述数据构成的群胚：$X\setminus Z$ 上的一个 $G$-主丛，$X^{\wedge}_Z$ 上的一个 $G$-主丛，以及二者在 $(X^{\wedge}_Z\setminus Z)$ 上的一个等同。
\end{lemma}

\begin{remark}
这里正是真正有用到 Fargues–Fontaine 曲线的代数版本的地方。
\end{remark}

\begin{example}
对 $G=\mathrm{GL}_2$，$\mu=(1,0)$，限制到 $\Gr^{B^+_{\mathrm{dR}}}_{\mathrm{GL}_2,\le\mu}\cong(\PP^1_E)^{\diamond}$ 就恢复了先前的例子。
\end{example}

\begin{theorem}
$\Gr^{B^+_{\mathrm{dR}}}_G\to\Bun_G$ 是 $v$-叠的满射。
\end{theorem}

\begin{remark}
这是 Drinfeld–Simpson 关于通常曲线的一个结果的类比。在那里只有当 $G$ 半单时才成立，对环面不成立。在此情形，它对所有约化 $G$ 都真的成立。关键在于去心曲线的 Picard 群平凡。（这显然是该陈述成立所必需的。）
\end{remark}

\begin{proof}
关键步骤是在几何点上检验会发生什么。这归功于 Fargues / Anschütz，用到 $G$-主丛的分类。

由此我们逐步推广到一般情形。对一般的 $S$，$\Bun_G$ 的 $S$-点对应 $\cE/X_S$。设 $S$ 严格完全间断。在几何点上可由上一步提升到 $\Gr^{B^+_{\mathrm{dR}}}_G$，得到 $\cE_s$ 的修正 $\cE'_s$。选取 $\cE$ 的一个修正 $\cE'$，使其在 $s$ 处恢复 $\cE'_s$。（用 Cartan 分解，这是提升矩阵的问题。）只需证明 $\cE'$ 在 $s$ 的某个邻域内平凡。我们断言，一般地对 $\cE'/X_S$，使 $\cE'|_{X_s}$ 平凡的轨迹在 $S$ 中是开的。这将在下次补完。
\end{proof}

\begin{remark}[Witt 向量仿射 Grassmann 簇]
设 $G/\cO_E$。Witt 向量仿射 Grassmann 簇 $\Gr^{\mathrm{Witt}}_G$ 是 $\Fq$ 上完满代数上的函子，取
\[
R\mapsto G(W_{\cO_E}(R)[1/\pi])/G(W_{\cO_E}(R)).
\]
它可由一个 ind(完满概形) 表示。

存在 $B^+_{\mathrm{dR}}$-仿射 Grassmann 簇的退化。我们可以定义 $\Gr^{B^+_{\mathrm{dR}}}_G\to\Spa\cO_E^{\diamond}$。这是 $\Spa\cO_E^{\diamond}$ 上的 ind-(紧合、在空间菱形中相对可表) 对象。这是从 $\Gr^{B^+_{\mathrm{dR}}}_G$ 到 $(\Gr^{\mathrm{Witt}}_G)^{\diamond}$ 的退化。
\end{remark}
